%!TEX root=../main.tex
\section{Introduction} % (fold)
\label{sec:introduction}


Due to the increasing availability of detailed performance data, sport analytics is a growing field in recent times and receives a large interest in the statistics community. Concerning team sports, traditional approaches concern the prediction of game results based on team characteristics \citep{dixonModellingAssociationFootball1997,babootaPredictiveAnalysisModelling2019}. See \cite{horvatUseMachineLearning2020,bunkerApplicationMachineLearning2022} for recent reviews. Recently, the availability of tracking data allowed researchers to focus on more specific aspects of different sports, such as evaluating individual player performance in soccer \citep{pappalardoPlayeRankDatadrivenPerformance2019,decroosActionsSpeakLouder2019} or defensive performance on passing plays in American football \citep{schmidSimulatingDefensiveTrajectories2021}. Because of the data provided by the National Basket Association (NBA), basketball is an interesting sport to analyse from this perspective.

Basketball is a sport where two teams, consising of five players, compete each others. The objective of the game is to shoot a ball through the opponent's hoop while preventing the ball to go through your own hoop. Basketball has been widely studied from different perspective. Shooting movements have been analyzed in \cite{millerRelationshipBasketballShooting1996,liuChangesBasketballShooting1999,okazakiReviewBasketballJump2015}. Estimating game win probability, in-game win probability \citep{maddoxBayesianEstimationIngame2022} or possession outcomes \citep{cervoneMultiresolutionStochasticProcess2016}. Defense analysis \citep{csataljayEffectsDefensivePressure2013}. Spatial performance \citep{zuccolottoSpatialPerformanceAnalysis2023}. See \cite{ternerModelingPlayerTeam2021} for a review on performance analysis in basketball.

In basketball, the outcome of the game is greatly determined by the shooting accuracy. An important factor of a player's effectiveness is its ability to consistently make successful shots from various locations on the court. The analysis of shooting patterns can thus evaluate individual player performance but also help to develop team strategies. An important tool to analyse shooting patterns is shooting charts. Shooting charts are graphic representations of all the shots a player has taken and made across a season or a game. Shooting charts have been widely analysed in the basketball litterature. Most of the work in this area concern the modelling of the shooting charts as point process. To do so, the idea is to construct a count matrix of the number of shots by player on a discretized court, fit an intensity surface for each each player over the discretized court using a log-Gaussian Cox process and find a low-rank approximation of all players' intensity surface using non-negative matrix factorization \citep{millerFactorizedPointProcess2014,franksCharacterizingSpatialStructure2015}. Non-parametric estimation of the intensity surface is also considered in \cite{yinBayesianNonparametricLearning2022}. The low-rank approximation of the intensity surfaces is then used to summarize the shooting habits of NBA players as well as characterizing their efficiency. Some authors developed models that are not based on fitting an intensity surface for each player. For examples, \cite{reichSpatialAnalysisBasketball2006} proposed conditionally autoregressive models to consider spatial correlation and \cite{huZeroinflatedPoissonModel2023} fit a Bayesian zero-inflated Poisson model to model shoot attempts of players with different shooting preference. These methods allow us to summarize information about players but not compare different players. 

To compare players regarding their shooting behavior, specific methods have been developed. \cite{fuHoopInSightAnalyzingComparing2024} use (interactive) data visualization to analyze and compare basketball shooting performance.  \cite{huBayesianGroupLearning2021} use a log-Gaussian Cox process and define a similarity measure to compare the players. \cite{jiaoBayesianMarkedSpatial2021} proposes a Bayesian marked spatial point processes model to model the shots intensity. They study if the sucess rate of shot is higher at locations where players makes more shots and provide a clustering of players. \cite{wong-toiJointAnalysisField2022} analyzes shooting patterns among different players. They quantify the shooting patterns using a discretisation of the field by performing a Bayesian nonparametric mixture model to find latent clusters of players. \cite{yinAnalysisProfessionalBasketball2023} proposes a Bayesian nonparametric matrix clustering approach on estimated intensity maps to analyze the structure in the shot selection. These approachs require a discretization of the court. 

To suppress the need of the discretization of the court, we propose to estimate the density of the shots continuously on the offensive half-court and to model these densities as functional data. We consider both attempted and made shots modeled as multivariate functional data. Functional data analysis concerns the analysis of realisations of a stochastic process, recorded at discrete times and defined on compact intervals. Considering the density of the shoots as a stochastic process defined over the rectangle of the offensive half-court, we view the denisty of the shoots of any individual player as a realisation of this process. We aim to the variability of this stochastic process. To do so, we perform dimension reduction using multivariate functional principal components analysis. We can then characterise the shooting style of the players using the eigencomponents. Finally, we can determine which players are similar to each other using the projection of the data into the basis defined by the eigencomponents and clustering techniques. By providing a more nuanced understanding of shooting patterns, we aim to offer valuable insights for coaches, analysts, and players.

The rest of the article is organized as follows. In Section~\ref{sec:dataset}, we introduce the motivating dataset from the 2018-2019 regular season to the 2023-2024 regular season. In Section~\ref{sec:method}, we present our methodology. Application of the methodology to NBA player data are reported in Section~\ref{sec:results}. We conclude the article in Section~\ref{sec:discussion} with a discussion.


% section introduction (end)